\section{集体做了一个错误决定：复杂度、产品与人}

第二章标题是“哦，我们集体做了一个错误决定！”这不是单个小失误，而是创业团队在技术、产品和组织之间反复试错的缩影。Peak 多次谈到复杂度：技术复杂、产品复杂、人更复杂。相比程序的复杂度，人和组织的复杂度更难控制。

\subsection{为什么人比 AI 更复杂}

Peak 说自己不喜欢管人，因为人的复杂度会随着人数增加近似指数增长。程序如果设计模式足够好，复杂度还可控；人和组织则会产生沟通、动机、期待、情绪和权责问题。这对 Agent 公司尤其重要：AI 可能提升个人生产力，但不自动解决组织复杂性。

\begin{warningbox}{AI 工具不会自动消除组织问题}
即使每个员工都用了 AI 工具，组织仍然需要目标、协作、决策、责任和信息流。AI 可以放大个人能力，也可能放大混乱。
\end{warningbox}

\subsection{AI 时代的组织变化}

前面讲人和组织的复杂性，本节看 AI 会怎样进入组织。Peak 认为 AI 时代组织变化没有大家想象那么大，但组织里会多出很多 AI。Manus 鼓励员工使用各种 AI 产品，甚至第三方工具也尽量报销，因为 AI 公司首先要让员工理解业界前沿。这个观点很实际：不是先重构组织图，而是让个体先获得两倍、十倍生产力。

\begin{importantbox}{组织变化的顺序}
先让个体理解并使用 AI，再观察工作流和组织结构如何变化。过早设计“AI 时代组织形态”，可能比让团队真实使用 AI 更慢。
\end{importantbox}

\subsection{本章小结}

Manus 的错误和试错，不只来自产品方向，也来自复杂系统中的人和组织。Agent 公司要处理的不只是模型能力，还有人如何协作、如何保持简单、如何不被功能和组织复杂度吞没。

